\section{Japan and the United Kingdom}\label{sec:international}

The United Kingdom and Japan also borrow in their own currencies, and both ran QE on a scale comparable to or larger than the Fed's. Each differs from the United States in a way that matters. The UK issues the longest debt among the major advanced economies, and a quarter of it is index-linked. Japan has the highest debt, has had $r$ below $g$ throughout our sample, and until March 2024 paid a zero rate on much of its reserves. We apply the U.S. definitions to both. All inflation-layer ratios in this section use $g=4\%$ and $H=10$.

\subsection{Data}\label{sec:intldata}

For the UK we use gilts in issue by ISIN at each year-end, 2007--2025, from the Debt Management Office; the sums match its stated totals in every year. Asset Purchase Facility holdings are rebuilt gilt by gilt from the Bank of England's operation-level results, and the current stock matches the Bank's table exactly. Reserves earn Bank Rate, and notes and coin are the zero-interest base; Treasury bills, 3--4\% of marketable debt, are not yet included. For Japan we use JGBs by issue from the Ministry of Finance's debt yearbook, fiscal years 2020--2024, and Bank of Japan holdings by issue. BoJ current accounts are split by the rate they earn: the 0\% tier of 2016 to March 2024, and required reserves after the reform, join banknotes in the zero-interest base. Appendix~\ref{app:data} gives the details.

For the fiscal threshold, $r$ prices each security at the year's yield for its original tenor, and $g$ is trailing ten-year real growth plus the 2\% inflation target the three central banks share. For the United States this gives $g=4.4\%$ in 2025, above the FOMC-based 3.8\% of Section~\ref{sec:pricing}, so comparisons in this section are within this common basis.

\subsection{The clock and the inflation layer}\label{sec:intlclock}

\begin{table}[htbp]
\centering
\caption{Three sovereigns, own and consolidated}\label{tab:threesov}
{\small g = 4\%, H = 10; Japan: fiscal years ending in March of the following year\par}\smallskip
\footnotesize
\begin{tabularx}{\textwidth}{>{\raggedright\arraybackslash}p{0.3\textwidth}>{\centering\arraybackslash}X>{\centering\arraybackslash}X>{\centering\arraybackslash}X>{\centering\arraybackslash}X>{\centering\arraybackslash}X>{\centering\arraybackslash}X>{\centering\arraybackslash}X}
\toprule
 & U.S. 2007 & U.S. 2025 & UK 2007 & UK 2021 & UK 2025 & Japan FY2022 & Japan FY2024 \\
\midrule
Consolidated interest-bearing debt, \% of GDP & 27 & 95 & 32 & 100 & 99 & 150 & 183 \\
Overnight share of consolidated debt, \% & 1 & 14 & 5 & 42 & 22 & 28 & 44 \\
Zero-interest base / interest-bearing debt, \% & 22 & 8 & 10 & 4 & 3.5 & 48 & 11 \\
Index-linked share of consolidated debt, \% & 11 & 6 & 27 & 22 & 23 & 1 & 0.5 \\
Average maturity, years: own / consolidated & 4.6 / 4.7 & 5.8 / 4.8 & 14.3 / 13.6 & 15.0 / 9.4 & 13.9 / 11.1 & 8.3 / 7.1 & 8.6 / 6.0 \\
P(1): own / consolidated & 0.38 / 0.37 & 0.37 / 0.46 & 0.09 / 0.13 & 0.10 / 0.48 & 0.10 / 0.28 & 0.26 / 0.51 & 0.24 / 0.59 \\
Inflation-layer ratio $\int P/\int E$: own / consolidated & 2.73 / 1.45 & 2.22 / 2.18 & 0.90 / 0.81 & 0.87 / 2.47 & 0.89 / 1.41 & 1.20 / 0.81 & 1.14 / 1.81 \\
\bottomrule
\end{tabularx}
\par\smallskip{\footnotesize\raggedright \emph{Notes.} ``Own'' is the government's own marketable debt, with no central bank. Consolidated debt nets central-bank and government holdings and adds interest-bearing reserves (and, for the United States, reverse repos). The zero-interest base is currency (banknotes, notes and coin) plus reserves that earn nothing: U.S. reserves before October 2008, and Japan's 0\% tier and required reserves.\par}
\end{table}


\begin{figure}[htbp]
\centering
\includegraphics[width=0.95\textwidth]{fig5_three_countries.pdf}
\caption{Three sovereigns}\label{fig:three}
\begin{minipage}{0.95\textwidth}\footnotesize United States (blue), United Kingdom (orange), Japan (green). A: inflation-layer ratio, consolidated (solid) and own debt (dashed), $g=4\%$, $H=10$. B: fiscal threshold $\varphi^*$ against the debt-sensitivity of rates $\psi$, latest year, $g$ = trailing ten-year real growth $+$ 2\%; the dotted line marks the U.S. baseline $\psi=3$bp. C: cumulative real transfer from the post-2020 inflation surprises, from the end-2020 consolidated balance sheet (Japan: end-March 2022, plotted at fiscal year-ends), in \% of GDP in the base year, under full Fisher repricing (B, dashed), inflation as priced (D, dotted; not available for Japan) and actual yields (C, solid).\end{minipage}
\end{figure}

Table~\ref{tab:threesov} and Figure~\ref{fig:three}A compare the three balance sheets. On their own debt, the three governments look very different. The UK's gilts average 14--16 years, and only a tenth of a rate shock reaches the average rate within a year; U.S. debt reprices fastest. With long and partly index-linked debt, the UK's own-debt ratio is about 0.9, against 1.1--1.3 for Japan and 2.1--2.7 for the United States. Judged by debt management alone, the UK is the best placed to inflate.

Consolidation erases most of the difference. At the QE peak the consolidated ratios were almost equal: 2.47 for the UK and 2.46 for the United States at end-2021, and 1.99 for Japan at end-March 2024. At end-2021 the APF held 36\% of gilts, financed by reserves paid Bank Rate. As a result, the UK's consolidated maturity fell from 15.0 to 9.4 years, and 42\% of its consolidated liabilities paid the overnight rate.

What matters is how the central bank funds its bonds. Before QE, consolidation lowered the ratio in the United States (2.73 to 1.45) and in the UK (0.90 to 0.81), because central-bank liabilities were mostly currency. Japan shows the mechanism at scale. In FY2022 the BoJ held 47\% of JGBs, but much of its current accounts sat in the 0\% tier; with banknotes, the zero-interest base was 48\% of interest-bearing consolidated debt, and the consolidated ratio, 0.81, was below Japan's own-debt ratio. The March 2024 reform moved all but required reserves to the policy rate. With almost no change in the government's own debt, the ratio rose to 1.99 and interest-bearing consolidated debt rose from 150\% to 194\% of GDP. This is the counterfactual of Section~\ref{sec:policy} run in reverse, and the rise, 1.18, exceeds the entire rise in the U.S. ratio from 2007 to 2021.

QT unwinds the effect. Sales and redemptions took the APF from 36\% of gilts in 2021 to 17\% in 2025. Over the same period the UK ratio fell from 2.47 to 1.41, while the own-debt clock did not move. The Fed, which ran off its holdings without sales, ended 2025 at 2.18.

The same funding sets the monetary tolerance threshold. Over ten years, $\kappa^*$ was 1.23 in the UK in 2007, 0.41 at the 2021 peak and 0.71 after QT; Japan's fell from 1.1--1.2 under the three-tier system to 0.50--0.55 after the reform. Over five years the levels are higher (2.39, 0.57 and 1.18 for the UK; 1.4--1.6 and 0.65--0.73 for Japan), but the pattern is the same: funding central-bank bonds with interest-bearing reserves cut the threshold sharply.

\subsection{The fiscal threshold}\label{sec:intlthreshold}

\begin{table}[htbp]
\centering
\caption{Fiscal threshold and inflation requirement on a common basis}\label{tab:threshold3}
{\small +1pp permanent rate rise, H = 10\par}\smallskip
\small
\begin{tabularx}{\textwidth}{>{\raggedright\arraybackslash}p{0.3\textwidth}>{\centering\arraybackslash}X>{\centering\arraybackslash}X>{\centering\arraybackslash}X}
\toprule
 & U.S. 2025 & UK 2025 & Japan FY2024 \\
\midrule
Marginal cost r / growth g, \% & 4.2 / 4.4 & 4.8 / 3.4 & 1.4 / 2.6 \\
r $-$ g, pp & $-$0.2 & +1.3 & $-$1.2 \\
Consolidated debt b, \% of GDP & 95 & 99 & 183 \\
$\varphi^*$ at $\psi$ = 0 / 1 / 3 / 4.5 bp & $-$0.05 / 0.14 / 0.37 / 0.48 & 0.28 / 0.41 / 0.56 / 0.63 & $-$0.90 / 0.18 / 0.62 / 0.73 \\
$\varphi^*$ at $\psi$ = 3bp, trailing nominal g & 0.24 & 0.39 & 0.72 \\
Required inflation, pp a year, $\psi$ = 3bp: $\hat\varphi$ = 0 / 0.25 & 0.82 / 0.27 & 0.79 / 0.44 & 1.12 / 0.66 \\
Required inflation, pp a year, $\psi$ = 0: $\hat\varphi$ = 0 / 0.25 & 0 / 0 & 0.40 / 0.05 & 0 / 0 \\
\bottomrule
\end{tabularx}
\end{table}


Figure~\ref{fig:three}B plots $\varphi^*$ against $\psi$, and Table~\ref{tab:threshold3} reports the values. The UK is the only one of the three with $r$ above $g$, by 1.3pp in 2025. Its threshold is therefore positive even if rates do not respond to debt, 0.28 at $\psi=0$, and with $\psi=0$ and no fiscal response it is the only one of the three that would need the inflation route at all. Japan's threshold depends almost entirely on $\psi$: $-0.90$ at $\psi=0$ and 0.62 at $\psi=3$bp. The U.S. debt sensitivity cannot be assumed for Japan, where yields stayed near zero while debt rose past twice GDP, and Japan's own history of yield control does not identify it. The United States sits on the boundary. At $r\approx g$ its threshold is close to zero without a debt premium; it is 0.37 at $\psi=3$bp on this basis, and 0.47 with the FOMC growth path of Section~\ref{sec:pricing}.

At $\psi=3$bp and no fiscal response, the requirement is about 0.8pp a year in the United States and the UK and 1.1pp in Japan. The UK has the higher threshold but the lower ratio; Japan has the highest threshold and, since 2024, a ratio close to the others.

\subsection{The 2021--25 inflation surprise in the United Kingdom}\label{sec:uktest}

We repeat the test of Section~\ref{sec:inflationtest} on the UK's end-2020 consolidated balance sheet: privately held conventional gilts (\pounds871bn), reserves (\pounds770bn) and notes and coin (\pounds92bn). Index-linked gilts (\pounds454bn) cannot be eroded and are reported separately. Expected CPI inflation is 2.39\%: the end-2020 five-year implied RPI inflation less the 2011--20 average RPI--CPI wedge. The design differs from the U.S. test in one respect: the portfolio is closed. Each gilt is rolled at redemption into one of the same original tenor, and amounts stay at their end-2020 levels, so new borrowing, QT and Treasury bills are left out. Yields are the Bank of England's zero-coupon curves, with Bank Rate at the short end.

\begin{table}[htbp]
\centering
\caption{United Kingdom: real transfer to the consolidated government from the 2021–25 inflation surprise}\label{tab:uktest}
{\small cumulative, \% of 2020 GDP\par}\smallskip
\footnotesize
\begin{tabularx}{\textwidth}{l>{\centering\arraybackslash}X>{\centering\arraybackslash}X>{\centering\arraybackslash}X>{\centering\arraybackslash}X>{\centering\arraybackslash}X>{\centering\arraybackslash}X>{\centering\arraybackslash}X>{\centering\arraybackslash}X}
\toprule
End of & Cumulative surprise (pp) & Gross erosion & Model: full Fisher repricing (B) & Model, analytic s$\cdot$b$\cdot$E & Inflation as priced (D) & Actual (C) & Memo: extra interest on reserves in C & Memo: RPI surprise on index-linked gilts \\
\midrule
2021 & 2.9 & 2.3 & 1.3 & 1.2 & 2.2 & 2.3 & 0.0 & 0.9 \\
2022 & 10.5 & 8.5 & 4.6 & 4.2 & 8.1 & 8.0 & 0.5 & 2.9 \\
2023 & 11.9 & 9.7 & 5.2 & 4.7 & 9.2 & 7.3 & 2.2 & 3.3 \\
2024 & 12.1 & 9.9 & 5.1 & 4.8 & 9.2 & 5.4 & 4.1 & 3.4 \\
2025 & 13.0 & 10.6 & 5.3 & 5.0 & 9.9 & 4.2 & 5.7 & 3.6 \\
\bottomrule
\end{tabularx}
\par\smallskip{\footnotesize\raggedright \emph{Notes.} Extra interest on reserves is measured relative to the no-surprise forward path. The RPI surprise on index-linked gilts is their inflation uplift in excess of the end-2020 implied RPI inflation. It is a nominal cost, not a real transfer, because holders of those gilts are fully compensated.\par}
\end{table}


The UK results repeat the U.S. pattern closely (Table~\ref{tab:uktest}, Figure~\ref{fig:three}C). Markets repriced far less than full Fisher repricing: with inflation as priced (D), the transfer reached 9.2\% of GDP by 2023, 1.8 times the full-Fisher 5.2\% (B), the same multiple as in the United States. Higher real rates took most of it back, almost all through reserves. With actual yields (C), the transfer peaked at 8.0\% of GDP in 2022 and fell to 4.2\% by 2025, the U.S. end-point. The 2025 gap between C and D, 5.7\% of GDP, is about the extra interest paid on reserves as Bank Rate rose to 5.25\%; reserves were 47\% of the UK's interest-bearing nominal liabilities at end-2020. Index-linked gilts shifted part of the inflation to the government. RPI exceeded its implied rate by 16.6pp over 2021--25, which added 3.6\% of GDP to nominal debt, about a third of the gross erosion. Debt that inflation cannot erode also cannot share in a surprise.

\subsection{The 2022--25 inflation surprise in Japan}\label{sec:jptest}

Japan's inflation came later and was lower, but it started from near-zero expectations, with yields held down by yield-curve control. At end-March 2022 the consolidated balance sheet held, in trillion yen, privately held nominal JGBs and bills of 659, interest-bearing current accounts of 262, current accounts in the 0\% tier of 301, and banknotes of 120: together 233\% of FY2021 GDP. No public breakeven series is available, so expected inflation is the ten-year average CPI inflation to March 2022, 0.61\%, and there is no run D. The closed-portfolio design follows the UK, over the 48 months to March 2026. Floating-rate JGBs reset every six months at the ten-year yield. Interest-bearing current accounts earn their tier-weighted rate until the March 2024 reform and the call rate after it; under actual yields (C), the former 0\% tier above required reserves also earns the call rate after the reform.

\begin{table}[htbp]
\centering
\caption{Japan: real transfer to the consolidated government from the 2022–25 inflation surprise}\label{tab:jptest}
{\small cumulative, \% of FY2021 GDP; fiscal years ending in March\par}\smallskip
\footnotesize
\begin{tabularx}{\textwidth}{l>{\centering\arraybackslash}X>{\centering\arraybackslash}X>{\centering\arraybackslash}X>{\centering\arraybackslash}X>{\centering\arraybackslash}X>{\centering\arraybackslash}X}
\toprule
End of & Cumulative surprise (pp) & Gross erosion & Model: full Fisher repricing (B) & Model, analytic s$\cdot$b$\cdot$E & Actual (C) & Memo: interest on the former 0\% tier in C \\
\midrule
FY2022 & 2.6 & 6.0 & 4.3 & 4.2 & 5.9 & 0.0 \\
FY2023 & 4.6 & 10.7 & 7.4 & 7.3 & 10.1 & 0.0 \\
FY2024 & 7.6 & 17.6 & 11.8 & 11.5 & 16.8 & 0.1 \\
FY2025 & 8.4 & 19.5 & 12.7 & 12.6 & 18.0 & 0.4 \\
\bottomrule
\end{tabularx}
\par\smallskip{\footnotesize\raggedright \emph{Notes.} With expected inflation of 1\% instead of 0.61\%, the actual transfer (C) at end-FY2025 is 14.5\% of GDP and the full-Fisher transfer (B) 10.4\%.\par}
\end{table}


The result differs from the U.S. and UK tests in the direction that matters (Table~\ref{tab:jptest}). The transfer was not taken back. Under actual yields it rose every year, to 18.0\% of FY2021 GDP by March 2026, 92\% of gross erosion and 1.4 times the full-Fisher benchmark. The ten-year yield rose from 0.2\% to 2.2\% while CPI inflation averaged about 3\% in FY2022--24, so real rates on new debt fell. The zero-interest base did much of the work: banknotes and the 0\% tier were 31\% of the eroded liabilities, and they never reprice. The 2024 reform cost little over this horizon (0.4\% of GDP), because the policy rate stayed low; its effect is on the price going forward. Japan thus shows the inflation route at its most effective, and why it is not repeatable. The conditions were a large zero-interest base, slow repricing, and yields that did not price the surprise. Since the reform, current accounts earn the policy rate and the consolidated ratio is about 2, so the same route would now be costlier.

\subsection{What the comparison adds}\label{sec:intlsummary}

Four points generalize beyond the United States. The inflation layer is a property of the consolidated balance sheet, not of debt management. How the central bank funds its bonds is the decisive margin: it sets both the price of the inflation route and the monetary tightening the route can survive. Which route is taken depends on the central bank: markets underpriced the surprise everywhere, and higher real rates then took the transfer back in the United States and the UK, but not in Japan. And the fiscal threshold is where the countries differ, through $r-g$ and the debt sensitivity of rates. Three sovereigns are not a panel, but they show that the variation in maturity, indexation and central-bank funding needed to identify the fiscal response across countries exists.
